\documentclass[12pt,a4paper]{ctexart}

\usepackage[margin=2.2cm]{geometry}
\usepackage{amsmath,amssymb,mathtools}
\usepackage{graphicx}
\usepackage{booktabs,longtable,array,tabularx}
\usepackage{enumitem}
\usepackage{xcolor}
\usepackage[most]{tcolorbox}
\usepackage{hyperref}
\usepackage{tikz}

\hypersetup{
  colorlinks=true,
  linkcolor=blue!55!black,
  urlcolor=blue!55!black
}

\setlist{itemsep=2pt,topsep=4pt}
\setlength{\parindent}{2em}
\setlength{\parskip}{3pt}

\newtcolorbox{questionbox}[1]{
  enhanced,
  breakable,
  colback=blue!3!white,
  colframe=blue!65!black,
  colbacktitle=blue!65!black,
  coltitle=white,
  title={#1},
  fonttitle=\bfseries,
  left=3mm,
  right=3mm,
  top=1.5mm,
  bottom=1.5mm,
  boxrule=0.75pt,
  arc=1mm
}

\newtcolorbox{answerbox}[1][正式答案]{
  enhanced,
  breakable,
  colback=green!4!white,
  colframe=green!45!black,
  colbacktitle=green!45!black,
  coltitle=white,
  title={#1},
  fonttitle=\bfseries,
  left=3mm,
  right=3mm,
  top=1.5mm,
  bottom=1.5mm,
  boxrule=0.75pt,
  arc=1mm
}

\newtcolorbox{notebox}[1][重点]{
  enhanced,
  breakable,
  colback=yellow!8!white,
  colframe=yellow!45!black,
  title={#1},
  fonttitle=\bfseries,
  left=3mm,
  right=3mm,
  top=1.5mm,
  bottom=1.5mm,
  boxrule=0.6pt,
  arc=1mm
}

\newtcolorbox{warningbox}[1][易错点]{
  enhanced,
  breakable,
  colback=red!3!white,
  colframe=red!55!black,
  title={#1},
  fonttitle=\bfseries,
  left=3mm,
  right=3mm,
  top=1.5mm,
  bottom=1.5mm,
  boxrule=0.6pt,
  arc=1mm
}

\newcommand{\examfield}[1]{\par\smallskip\noindent\textbf{#1：}}
\newcommand{\sourcehint}[1]{\par\smallskip\noindent{\footnotesize\color{gray!75!black}来源：#1}}

\title{课程期末复习资料}
\author{}
\date{\today}

\begin{document}

\maketitle
\tableofcontents
\newpage

\section{考试范围与复习总览}

% 根据真实资料填写课程框架、考试范围、教师重点、章节关系和复习顺序。

\section{理论课章节知识点完整总结}

% 大课程建议拆成 sections/ch01-title.tex 等文件，并在此处使用 \input{sections/ch01-title} 组装。
% 课件例题应放在对应知识点附近；作业题、往年题、专项训练和预测题放在后续独立章节。

\subsection{示例章节标题}

\begin{notebox}[本章目标]
这里写本章要掌握的核心问题。正式资料中删除所有示例文字。
\end{notebox}

\subsubsection{核心知识点}

写成完整讲解：定义、背景、原理、公式、适用条件、例题、图表解释、易错点和考试问法。

\subsubsection{课件例题}

把教师课件中的例题放在对应知识点下，先说明它考查什么，再写完整解析和答案。

\section{公式与关键结论汇总}

\section{易混概念、方法与模型对比}

\section{作业题与概念理解整理}

\subsection*{作业题 X：题目标题}

\begin{questionbox}{原题}
放原题题干。题干和答案必须分开。
\end{questionbox}

\examfield{对应知识点}
填写知识点。

\examfield{题目解析}
填写解析。

\begin{answerbox}
填写正式答案。
\end{answerbox}

\begin{warningbox}
填写易错点。
\end{warningbox}

\section{往年试卷整理}

\section{考试题型专项训练与预测题}

% 根据用户提供的考试题型组织训练，例如名词解释、简答、选择、判断、计算、应用、代码分析或设计题。
% 答题思路只从教师答案、作业答案、样卷答案或往年参考答案中归纳，不硬造固定模板。
% 若考试会考名词解释或术语解释，可另建“名词解释集中背诵版”章节或单独输出一份背诵版。
% 背诵版条目保持干净：术语、英文名（必要时）、解释、辨析或例子；不要加入 [原题][核心] 等标签。

\section{考试当天速览}

% 自由组织成适合临考快速查看的形式，例如关键概念、易错点、公式、题型步骤或最后检查清单。

\end{document}
